\clearpage
\thispagestyle{empty}
\newgeometry{left=20mm, right=20mm, top=20mm, bottom=20mm}

\begingroup
\sffamily
\footnotesize
\setlength{\parindent}{0pt}
\setlength{\parskip}{0.6em plus 0.2em minus 0.1em}

\noindent
\textbf{УДК}~517(075.8)\\
\textbf{ББК}~22.161я73\\
\phantom{\textbf{ББК}~}Ч-50

\vspace{1.4em}

{\bfseries X.~Y.~Zettov, И.~И.~Иванов}

\vspace{0.4em}

{\bfseries \LaTeX-ШАБЛОН УЧЕБНИКА. Часть~I.}~--- Москва, 2026.~--- \pageref*{LastPage}~с.

\vspace{1.4em}

Настоящее издание представляет собой демонстрационный документ и одновременно готовый шаблон \LaTeX-проекта для подготовки учебных и научно-технических материалов по математике, информатике и смежным дисциплинам.
Шаблон может использоваться как отправная точка для создания конспектов лекций, учебных пособий, методических материалов, курсовых и дипломных работ, диссертаций и полноформатных книг.

В шаблоне реализована единая система семантических блоков для определений, утверждений, лемм, теорем, следствий, примеров, замечаний, алгоритмов, программного кода и задач.
Поддерживаются общая нумерация блоков, именованные гиперссылки, ненумерованные блоки, автоматическое оформление доказательств и решений, а также специальные маркеры дополнительного материала и необязательных доказательств.
Отдельное внимание уделено математической типографике, оформлению заголовков, формулам, рисункам, таблицам, спискам и единой системе обозначений.

Шаблон предусматривает вставку исходного кода на языке Python из внешних файлов, автоматическое отображение вывода программы и визуальное сопоставление строк кода со строками вывода.
Для записи вычислительных процедур подключён и русифицирован пакет \texttt{algpseudocode}.
В состав проекта также входит подробное руководство по стилю оформления, которое может использоваться авторами и редакторами как справочник, а большими языковыми моделями~--- как набор формальных требований к обработке исходного текста.

Документ собирается с помощью LuaLaTeX и сам служит демонстрацией большинства возможностей шаблона.
Исходный код шаблона распространяется свободно при условии указания автора дизайна и \LaTeX-оформления в итоговом документе.

\vfill

\begin{center}
    \includegraphics[height=16mm]{l-00/logo_airi.png}%
    \hspace{14mm}%
    \includegraphics[height=16mm]{l-00/logo_fai.png}%
\end{center}

\vspace{1.6em}

{\color{xgray!35}\hrule height 0.4pt}

\vspace{1.0em}

\noindent
\begin{tabular}{@{}r@{\hspace{0.6em}}l@{}}
    \textcopyright{} & X.~Y.~Zettov, И.~И.~Иванов, 2026 \\
\end{tabular}

\noindent
{\color{xgray}%
    \textcopyright{} Дизайн и \LaTeX-оформление: А.~В.~Чертков.\\
}

\endgroup
\restoregeometry